%% ====================================================================
%%  Chapter 7 — Future Work
%% ====================================================================
\chapter{Future Work}
\label{ch:futurework}

Beyond the results presented in this thesis, there are several directions for future work on this
project.

\paragraph{Deploying the model on E-JUST's edge devices.} The most immediate next step in this project is
to deploy the MobileNetV2-UNet model on edge devices on the E-JUST campus. While the computational cost
and accuracy of the model are reported in this thesis, the model has yet to be deployed on any device.
Because the inference time of the MobileNetV2-UNet on a CPU is 27.2\,ms compared to DLoc's 270.9\,ms, the
MobileNetV2-UNet is a strong candidate for deployment on edge devices.

\paragraph{Creating a new localization dataset on E-JUST's campus.} The model was trained and evaluated
on the WILD dataset, which exists in only one building: Jacobs Hall. A new dataset could be created at
E-JUST containing CSI measurements taken in a different building on the campus. Such a dataset would allow
the results of this thesis to be evaluated for their transferability to a new WiFi environment.

\paragraph{Statistical validation of the student models.} The statistical validation of the baseline
model is presented in this thesis. Validation of the student models under multiple random seeds was
de-scoped due to the computational cost of training those models. The student models could be trained
under several random seeds to provide significance testing of the comparisons between the students and the
baseline model.

\paragraph{Applying corruption augmentation during training.} The robustness of the models was tested by
applying corruptions only during the testing phase. If these corruptions were also applied during
training, it is likely that the robustness of the models would increase, which would be beneficial in
environments where such robustness is required.

\paragraph{Exploring other architectures.} The architecture of the MobileNetV2-UNet model was explored in
this research. Further exploration of other architectures---such as models that use Mamba layers or add
attention mechanisms to the skip connections---could reveal additional benefits to the student models.
